\chapter{谐振子}

\section{线性微分方程}

在学习物理学时，通常是把课程分成一系列的科目，如力学、电学、光学，等等，并且总是一门课程接着一门课程地学习。例如，到目前为止本门课程主要讨论的是力学。但是，有一件奇怪的事情却一再出现：即在物理学的不同领域中，甚至在其他的学科中，出现的方程式往往几乎是完全一样的，因此在这些不同领域中很多现象都有其类似之处。举一个最简单的例子，声波的传播在很多方面就与光波的传播相类似。如果我们深入地研究声学，就会发现要做的很多工作与我们深入研究光学时相同。所以，对一个领域中某种现象的研究可以扩展我们对另一个领域的知识。最好从一开始就认识到这种扩展是可能的，否则，人们就可能对为什么要花这么多的时间和精力来研究仅仅是力学中的很小一部分，感到不可理解。

我们将要学习的谐振子，在许多其他领域中都有相类似的东西，虽然我们从力学的例子，如挂在弹簧上的重物，小振幅的摆，或者某些其他的力学装置出发，但实际上我们是在学习某一种微分方程。这种方程在物理学和其他学科中反复出现，而且事实上它是许多现象中的一部分，是值得我们认真研究的。包含这个方程式的现象有：挂在弹簧上的一个具有质量的物体的振动；在电路中电荷的来回振荡；正在产生声波的音叉的振动；电子在原子中产生光波的类似振动；描写调节温度的恒温器之类的伺服系统的操作方程；化学反应中一些复杂的相互作用；在养料供给和细菌产生的毒素共同作用下菌落的繁殖和生长；狐狸吃兔子，兔子吃青草等等；所有这些现象遵循一些彼此非常相似的方程式，这就是为什么我们要这样详细地研究机械振子的原因。这些方程称为常系数线性微分方程。一个常系数线性微分方程包含几项之和，每一项都是因变量对自变量的微商再乘以某一个常数。如

